式\eqref{eq:q2-objective}--式\eqref{eq:q2-soc}及其边界均为线性关系，日前和日内子问题均以线性规划求解。为防止数值退化造成无意义的同时充、放电，在经济目标最优值的货币容差内再最小化电池吞吐量；该词典序次目标不改变购电成本。

\noindent\textbf{预测--计划--执行闭环算法}
\algorithmstart
\algostep{1}{构造历史信息}
在第 $d$ 日 $0{:}00$，只读取 $d$ 之前已经结束的负荷、光伏、价格和 SOC 记录，计算式\eqref{eq:q2-forecast}并更新完整残差池。
\algostep{2}{滚动选择参数}
若到达新的 $14$ 日校准块，则对所有 $(m,\alpha)$ 候选在此前验证窗口运行同一闭环政策，按实际成本选择候选；否则沿用当前已选参数。
\algostep{3}{生成日前额度}
对残差池执行 K-medoids，按式\eqref{eq:q2-reserve}计算风险储备，在基准预测上求解日前 LP，锁定 $q_{d,1:T}$。
\algostep{4}{逐时段因果执行}
每一时段仅读取当前实际负荷、光伏和上期 SOC；以已观测残差前缀更新未来预测，在剩余锁定额度下求解 MPC，并只实施当前一步的 $x,e,c,r,w$。
\algostep{5}{跨日传递与记账}
将 $E_{d,T}$ 传递至下一日，按 $\sum_t(p_tq_{d,t}+5p_te_{d,t})$ 写入普通计划费、紧急费和总成本台账。
\algostep{6}{审计与输出}
检查历史截止日、日内无未来信息、$x\le q$、能量平衡、SOC、功率及不同时充放电；随后按题设格式导出购电计划、充放电和紧急购电结果。
\algorithmend

考虑到 1 月历史样本不足，1 月只用于预测历史与 SOC 预热；正式题设工作簿导出 2 月 1 日至 12 月 31 日。该处理必须和完整年度核算分开报告，不能将两种口径混称为“全年结果”。
